\chapter{模型转换系统设计与实现}
\label{chap:5}
\chapterguide{交代研究方法怎样进入生产读写流程，突出公共数据与原生执行的边界，避免把各系统列成大量接口清单。}

\section{总体架构与模块分工}
\label{sec:5-1}

\subsection{公共模型层与系统适配层}
\label{subsec:5-1-1}
\outlinehint{说明CADExchange承担统一结构和公共服务，CADProj承担三系统读写、运行调度与展示，给出模块关系图。}

\subsection{原生读取、参数映射与创建执行的分离}
\label{subsec:5-1-2}
\outlinehint{以现有Creo组织说明快照、纯数值映射、写入计划与原生执行的分工，不将三端描述为已经完全相同的结构。}

\section{三系统的读取与重建适配}
\label{sec:5-2}

\subsection{SolidWorks特征读取与选择重建}
\label{subsec:5-2-1}
\outlinehint{概述特征遍历、参数读取、选择对象组织和创建，重点说明参考选择、不同生成分支及附属选项。}

\subsection{Creo元素结构与截面重建}
\label{subsec:5-2-2}
\outlinehint{说明特征参数与截面的读取、独立草图和内嵌截面的组织，以及引用、方向判定与目标创建计划。}

\subsection{NX对象构造与参数适配}
\label{subsec:5-2-3}
\outlinehint{说明对象和构造器路径、草图坐标反投影、对称深度和双向角度适配，以及辅助对象的依赖限制。}

\section{模型读取与顺序重建流程}
\label{sec:5-3}

\subsection{源模型历史、参数与引用提取}
\label{subsec:5-3-1}
\outlinehint{按历史读取五类特征、草图和引用，同时保留单位及必要几何，区分普通图元与构造图元。}

\subsection{目标特征创建与再生状态管理}
\label{subsec:5-3-2}
\outlinehint{按已具备的引用创建目标对象，取得建后实际坐标和可选几何；辅助对象按实际需要插入，不要求节点数完全相同。}

\subsection{运行配置、批处理与模型保存}
\label{subsec:5-3-3}
\outlinehint{交代模板、系统版本、工作单位、执行方式和输出路径；统计完整重建而非仅进程退出或文件生成。}

\section{结果记录、比较与可视化}
\label{sec:5-4}

\subsection{几何快照与源目标对应记录}
\label{subsec:5-4-1}
\outlinehint{说明逐特征几何和对象映射如何产生及复用，保留阶段、来源与状态，支持错误定位。}

\subsection{结构、参数与几何比较的接入}
\label{subsec:5-4-2}
\outlinehint{区分现有几何比较能力与待接回的目标重读、当前字段比较；介绍差异展示，不声称默认完整验证已闭合。}

\subsection{异常、受限处理与日志组织}
\label{subsec:5-4-3}
\outlinehint{列出未支持、读取失败、引用失败、创建失败、语义变化和未完成验证的记录需求，标注哪些仍需完善。}

\section{实现范围与能力边界}
\label{sec:5-5}

\subsection{五类特征的读写覆盖矩阵}
\label{subsec:5-5-1}
\outlinehint{按系统、特征、子类型、读取和创建分别列出范围，保留结构可表达与实际执行的差别。}

\subsection{辅助构造、受限参数与未完成部分}
\label{subsec:5-5-2}
\outlinehint{说明草图驱动尺寸、点与椭圆、特殊终止、不等量薄壁和复杂扫掠等限制，不将局部样例外推为完整支持。}

\section{本章小结}
\label{sec:5-6}

\outlinehint{归纳本章能够支持的认识，并说明与下一章的衔接；不重复列举所有接口或扩展未证明结论。}
